\section{\texorpdfstring{Filtered derived functors,\allowbreak{} page bounds,\allowbreak{} and the first Ext comparison}{Filtered derived  page  and the first Ext comparison}}\label{proofs-8352-8772-filtered-derived-functors-page-bounds-and-the-first-ext-comparison}

Editorial review of the Stacks Project authors\textquotesingle{} Derived Categories,\allowbreak{} original commit a04446e5\allowbreak{}7ec1fbc2\allowbreak{}52a871af\allowbreak{}cec7752f\allowbreak{}b2807b14, lines 8352–\allowbreak{}8772. SOURCE\_\allowbreak{}READING\_\allowbreak{}PART12.json records the frozen identities and the actual reading ranges. Earlier proof files remain part of this review,\allowbreak{} in particular PROOFS\_\allowbreak{}7096\_\allowbreak{}7578.md and PROOFS\_\allowbreak{}7589\_\allowbreak{}8350.md. The extensions below are editorial consequences; neither novelty nor source-body replacement is claimed.

\subsection{1. The set of objects needed for a calculation}\label{proofs-8352-8772-the-set-of-objects-needed-for-a-calculation}

Original 8352–\allowbreak{}8366 invokes current Sets,\allowbreak{} lemma-abelian-injectives,\allowbreak{} 1135–\allowbreak{}1157. That lemma includes more than an exact embedding of a small abelian category:\allowbreak{} it requires its injective objects to be exactly the objects that are injective in the ambient category. This is the property needed to compare the resulting resolutions.

Here is the construction in the usual locally small big-category setting of that statement. Start with the given set of objects and,\allowbreak{} when a set of diagrams is specified,\allowbreak{} every object occurring in those diagrams. Inductively enlarge a set S\_\allowbreak{}n of objects to a set S\_\allowbreak{}\{n\allowbreak{}+1\} containing:\allowbreak{} the zero object and chosen finite biproducts of objects of S\_\allowbreak{}n; chosen kernels and cokernels of every morphism between objects of S\_\allowbreak{}n; a chosen monomorphism X\allowbreak{}->J\_\allowbreak{}X with J\_\allowbreak{}X injective for every X in S\_\allowbreak{}n; and witnesses to noninjectivity for every X in S\_\allowbreak{}n that is not injective in the ambient category. Such a witness consists of a monomorphism U\_\allowbreak{}X\allowbreak{}->V\_\allowbreak{}X and a map U\_\allowbreak{}X\allowbreak{}->X admitting no extension V\_\allowbreak{}X\allowbreak{}->X; include U\_\allowbreak{}X and V\_\allowbreak{}X. Each stage has a set of requests:\allowbreak{} the category is locally small,\allowbreak{} there are a set of pairs of existing objects,\allowbreak{} and each witness involves finitely many objects and maps. Set-indexed collection and choice make the enlargement a set. This does not select data simultaneously for every object of the proper class.

Let A\textquotesingle{} be the full subcategory on the union of the S\_\allowbreak{}n. It is additive and has kernels and cokernels computed in A,\allowbreak{} because every finite collection of its objects appears at one stage. The comparison from coimage to image is the ambient isomorphism,\allowbreak{} so A\textquotesingle{} is abelian and its inclusion is exact. Every object X embeds in its J\_\allowbreak{}X within A',\allowbreak{} hence A\textquotesingle{} has enough injectives. An ambient injective in A\textquotesingle{} is injective there:\allowbreak{} all monomorphisms of A\textquotesingle{} remain monic in A,\allowbreak{} and fullness includes the ambient extensions. Conversely a noninjective X in A\textquotesingle{} has its chosen witness in a later stage of A',\allowbreak{} so is not injective in A\textquotesingle. This proves the additional injectivity property rather than assuming that exactness of an embedding alone would imply it.

For a set of filtered complexes,\allowbreak{} include all their terms,\allowbreak{} all their filtration subobjects,\allowbreak{} and the objects in their specified diagrams in S\_\allowbreak{}0. There are only set many:\allowbreak{} the degrees and filtration indices are integers. Each such diagram is then a diagram in A\textquotesingle. A filtered injective resolution formed there is also one in A,\allowbreak{} since its graded injectives are ambient injectives. The homotopy Hom groups between these complexes agree by fullness. Computing bounded-below filtered derived Hom by these resolutions therefore proves that the inclusion on these derived objects is fully faithful and computes the same maps. The analogous argument without filtrations proves the bounded-below unfiltered comparison. Given two such choices of A',\allowbreak{} apply the same construction to the union of their object sets; the resulting comparison uses the same ambient Hom groups. This proves compatibility of the calculations and does not assert a quasi-inverse selected on the entire proper class.

The remark\textquotesingle s ``only if'' refers to its preceding use of ordinary choice to pick representatives for a set of objects. It is not used here as an impossibility theorem about every big category. For example,\allowbreak{} PART11 supplies an explicit functorial construction on a fixed filtration window when functorial embeddings are given. No new source defect is counted from that reading of the remark. RECON-063 /\allowbreak{} DERIVED-058 reconciles the existing MC-STK-ERR-0869 possessive correction at 8364.

\subsection{2. A canonical functor on split filtered objects}\label{proofs-8352-8772-a-canonical-functor-on-split-filtered-objects}

Let T:\allowbreak{}A\allowbreak{}->B be additive. For a filtered injective I,\allowbreak{} all the sequences \[
 0\to F^{p+1}I\longrightarrow F^p I
      \longrightarrow \operatorname{gr}^p I\to0              \tag{2.1}
\] and all inclusions F\^{}p I\allowbreak{}->I split,\allowbreak{} by PART11 section 3. Thus T(F\^{}p I)\allowbreak{}\allowbreak{}->T(I)\allowbreak{} is monic. Give T(I)\allowbreak{} the filtration consisting of those actual subobjects. It is finite and is functorial:\allowbreak{} a filtered map f:\allowbreak{}I\allowbreak{}->J gives the commutative square from F\^{}p I\allowbreak{}->I to F\^{}p J\allowbreak{}->J,\allowbreak{} whose T-image restricts T(f)\allowbreak{} to the required subobjects. Addition,\allowbreak{} identity and composition are those of T.

If a decomposition I\allowbreak{}=direct-sum I\_\allowbreak{}p is chosen as in the source,\allowbreak{} the finite direct-sum comparison identifies this filtered T(I)\allowbreak{} with direct-sum T(I\_\allowbreak{}p)\allowbreak{},\allowbreak{} with filtration direct-sum\_\allowbreak{}\{p>\allowbreak{}=a\}T(I\_\allowbreak{}p)\allowbreak{}. A filtered map in such coordinates has components f\_\allowbreak{}\{q p\}:\allowbreak{}I\_\allowbreak{}p\allowbreak{}->J\_\allowbreak{}q with f\_\allowbreak{}\{q p\}\allowbreak{}=0 for q\textless p. Apply T to each component. Composition of the matrices is preserved because all relevant sums are finite and T is additive. A different splitting changes coordinates by an invertible filtration-preserving matrix; applying T gives an invertible change of coordinates with the same composition identities. This supplies the omitted map construction and proves independence from the noncanonical splittings in 8415–\allowbreak{}8439.

Applying T to (2.1)\allowbreak{} preserves its split exactness. The resulting quotient map is therefore a canonical natural isomorphism \[
 \operatorname{gr}^p(T_{\rm ext}I)
       \xrightarrow{\sim}T(\operatorname{gr}^p I).            \tag{2.2}
\] It is induced by the original quotient in (2.1)\allowbreak{},\allowbreak{} and does not require choosing a splitting. Forgetting the filtration similarly identifies T\_\allowbreak{}ext(I)\allowbreak{} with T(I)\allowbreak{},\allowbreak{} naturally. These assertions hold for every finite split filtered object,\allowbreak{} not only for injectives.

They do not assert commutation on arbitrary filtered objects. For example fix a prime ell,\allowbreak{} let T(M)\allowbreak{}\allowbreak{}=M[ell]\allowbreak{},\allowbreak{} and filter Z by F\^{}0 Z\allowbreak{}=Z,\allowbreak{} F\^{}1 Z\allowbreak{}=ell Z,\allowbreak{} F\^{}2 Z\allowbreak{}=0. Then T(Z)\allowbreak{}\allowbreak{}=0,\allowbreak{} so gr\textasciicircum{}0(T(Z)\allowbreak{})\allowbreak{}\allowbreak{}=0,\allowbreak{} whereas T(gr\^{}0 Z)\allowbreak{}\allowbreak{}=T(Z/\allowbreak{}ell Z)\allowbreak{}\allowbreak{}=Z/\allowbreak{}ell Z. This T is left exact:\allowbreak{} an element of the kernel of a map on ell-torsion is exactly the ell-torsion of its original kernel. Thus even left exactness does not make arbitrary associated graded objects commute with T. Original 8379 already gives that warning. The commutation mentioned at 8448 is the restricted one (2.2)\allowbreak{},\allowbreak{} in the context of T\_\allowbreak{}ext on filtered injectives. No extra defect is counted from that contextual shorthand.

\subsection{\texorpdfstring{3. DERIVED-UNDER-015:\allowbreak{} the image filtration before deriving}{3.  the image filtration before deriving}}\label{proofs-8352-8772-derived-under-015-the-image-filtration-before-deriving}

There is also an additive extension to all finite filtered objects,\allowbreak{} without left exactness. Keep the original T(A)\allowbreak{} and define \[
 T_{\rm im}(A,F)=
 \left(T(A),\
 F^pT_{\rm im}(A)=
 \operatorname{im}\big(T(F^pA)\longrightarrow T(A)\big)\right).
                                                               \tag{3.1}
\] The map in (3.1)\allowbreak{} is T of the actual filtration inclusion. The images decrease with p since the inclusions factor. If F\^{}a A\allowbreak{}=A and F\textasciicircum{}\{b\allowbreak{}+1\}A\allowbreak{}=0,\allowbreak{} the corresponding maps are an identity and a zero map,\allowbreak{} respectively. Hence (3.1)\allowbreak{} has the same permitted finite bounds [a,\allowbreak{}b]\allowbreak{}.

For a filtered map f:\allowbreak{}A\allowbreak{}->B the T-image of its filtration square maps the subobject in (3.1)\allowbreak{} into the corresponding subobject for B. This defines T\_\allowbreak{}im(f)\allowbreak{} by restricting the original map T(f)\allowbreak{}. It preserves identities and composition because T does. It is additive on Hom groups:\allowbreak{} the sum of two maps preserving these fixed subobjects preserves them,\allowbreak{} and its underlying map is T(f)\allowbreak{}\allowbreak{}+T(g)\allowbreak{}\allowbreak{}=T(f\allowbreak{}+g)\allowbreak{}. Images of the finite direct-sum inclusions show that finite biproducts and their filtrations are preserved. This proves that T\_\allowbreak{}im is an additive functor Fil\textasciicircum{}f(A)\allowbreak{}\allowbreak{}->Fil\textasciicircum{}f(B)\allowbreak{}. If T is left exact,\allowbreak{} the filtration maps T(F\^{}p A)\allowbreak{}\allowbreak{}->T(A)\allowbreak{} are monic,\allowbreak{} and (3.1)\allowbreak{} is exactly the source\textquotesingle s extension at 8375–\allowbreak{}8378. On split objects it is the canonical T\_\allowbreak{}ext of section 2 for every additive T.

Applying T\_\allowbreak{}im to complexes preserves differentials,\allowbreak{} filtered homotopies,\allowbreak{} shifts and cone matrices because it is additive. It consequently gives an exact functor on the filtered homotopy categories. Assume A has enough injectives and let u:\allowbreak{}K\allowbreak{}->I be a bounded-below filtered injective resolution. For any filtered quasi-isomorphism s:\allowbreak{}K\allowbreak{}->L,\allowbreak{} the filtered Hom comparison gives a unique homotopy class v:\allowbreak{}L\allowbreak{}->I with vs\allowbreak{}=u. The map v is a filtered quasi-isomorphism by two-out-of-three. Between filtered injective resolutions these comparison maps are homotopy equivalences,\allowbreak{} with inverse and homotopy identities obtained by the same uniqueness.

Thus the denominator diagram defining the right derived functor of T\_\allowbreak{}im has the same essential-constancy witness as in the injective-computation proof already established:\allowbreak{} refine to I,\allowbreak{} where every further injective comparison is invertible up to filtered homotopy. The actual derived value is \[
 R(T_{\rm im})(K)=T_{\rm im}(I)=T_{\rm ext}(I)
       \quad\text{in }DF^+(B),                              \tag{3.2}
\] with comparison induced by T\_\allowbreak{}im(u)\allowbreak{}. The witness proves existence and the canonical derived comparison; (3.2)\allowbreak{} is not merely a coincidence of object values. Its naturality on roofs follows from the comparison equation vs\allowbreak{}=u and its uniqueness. Hence it agrees with the source\textquotesingle s RT\allowbreak{}=T\_\allowbreak{}ext j',\allowbreak{} including the implicit localization from K\textasciicircum{}\allowbreak{}+(Fil\textasciicircum{}f B)\allowbreak{} to DF\textasciicircum{}\allowbreak{}+(B)\allowbreak{}.

This extension does not assume that T preserves monomorphisms or exact sequences in A. For a strict example use T(M)\allowbreak{}\allowbreak{}=M direct-sum (M/\allowbreak{}ell M)\allowbreak{} from PART10 section 2. Applied to the monomorphism ell:\allowbreak{}Z\allowbreak{}->Z,\allowbreak{} its second component is zero on Z/\allowbreak{}ell Z,\allowbreak{} so T is not left exact. Formula (3.1)\allowbreak{} nevertheless supplies the filtration by the actual images,\allowbreak{} including that lost component,\allowbreak{} and (3.2)\allowbreak{} supplies its derived functor. No equality with T(gr A)\allowbreak{} on an arbitrary underived filtration is asserted.

\subsection{\texorpdfstring{4. Derived associated graded objects,\allowbreak{} convergence,\allowbreak{} and naturality}{4. Derived associated graded   and naturality}}\label{proofs-8352-8772-derived-associated-graded-objects-convergence-and-naturality}

Let K\allowbreak{}->I be a filtered injective resolution as above. Every gr\^{}p I and the underlying I are bounded-below complexes of injectives,\allowbreak{} and the corresponding maps from gr\^{}p K and underlying K are quasi-isomorphisms. Formula (2.2)\allowbreak{} is natural in the maps of filtered objects,\allowbreak{} hence commutes with the original differentials of I. It gives the natural isomorphisms of derived functors \[
 \operatorname{gr}^p RT(K)\simeq RT(\operatorname{gr}^p K),
 \qquad
 \operatorname{forget}RT(K)\simeq RT(\operatorname{forget}K).
                                                               \tag{4.1}
\] The RT on the right is the unfiltered one. The naturality on arbitrary derived maps is checked on the injective comparison maps; inverting a denominator inverts that same commuting square. Consequently the statement is independent of j\textquotesingle{} and of chosen splittings. RECON-064 /\allowbreak{} DERIVED-061 is the existing MC-STK-ERR-0872 singular ``Lemma'' at 8447.

Put C\allowbreak{}=T\_\allowbreak{}ext(I)\allowbreak{},\allowbreak{} with its actual finite term filtrations. Its filtered-complex spectral sequence has \[
 E_0^{p,q}=\operatorname{gr}^p C^{p+q},\qquad
 E_1^{p,q}=H^{p+q}(\operatorname{gr}^p C)
          =R^{p+q}T(\operatorname{gr}^p K),                  \tag{4.2}
\] and differential of bidegree (r,\allowbreak{}1-r)\allowbreak{}. The first equality is the definition of the filtered spectral sequence,\allowbreak{} while the last is (2.2)\allowbreak{} applied to the injective resolution of gr\^{}p K. The boundedness,\allowbreak{} convergence and finite filtrations follow directly from the finite term filtrations,\allowbreak{} with the complete explicit calculation in section 5 below. The abutment is H\textasciicircum{}*(C)\allowbreak{}\allowbreak{}=R\textasciicircum{}*T(K)\allowbreak{}; each degree H\textasciicircum{}m(C)\allowbreak{} has its own finite filtration. RECON-065 /\allowbreak{} DERIVED-062 reconciles both MC-STK-ERR-0873 changes at 8530–\allowbreak{}8531. They state the graded abutment and its per-degree filtrations precisely; no different cohomology object is introduced.

For a filtered chain map K\allowbreak{}->L,\allowbreak{} the filtered Hom comparison provides a unique homotopy class I\allowbreak{}->J between their filtered injective resolutions,\allowbreak{} compatible with the augmentations. T\_\allowbreak{}ext,\allowbreak{} equivalently T\_\allowbreak{}im on these objects,\allowbreak{} preserves filtered homotopies. Such a homotopy induces a homotopy on the associated graded complexes,\allowbreak{} so equal maps on E\_\allowbreak{}1. Equality on each later page follows by passing to cohomology of the preceding page. The same uniqueness proves identity and composition. Thus the sequence is functorial and independent of choices from E\_\allowbreak{}1,\allowbreak{} exactly as stated at 8501. It need not have a choice-independent E\_\allowbreak{}0. For example,\allowbreak{} resolve Q[0]\allowbreak{} either by Q[0]\allowbreak{} or by Q[0]\allowbreak{} plus the contractible complex Q\allowbreak{}->Q in degrees 1,\allowbreak{}2,\allowbreak{} with a single filtration step. The E\_\allowbreak{}0 complexes differ whereas their E\_\allowbreak{}1 and later pages agree.

All of this argument uses additivity and split filtrations on I,\allowbreak{} and therefore also proves the additive version of the filtered spectral- sequence lemma. The source itself already announces additive RT at 8407. This propagates DERIVED-UNDER-011 through its actual filtered construction,\allowbreak{} with no suggestion of novelty.

\subsection{\texorpdfstring{5. DERIVED-UNDER-016:\allowbreak{} explicit local stabilization pages}{5.  explicit local stabilization pages}}\label{proofs-8352-8772-derived-under-016-explicit-local-stabilization-pages}

Let (C,\allowbreak{}F)\allowbreak{} be any complex in an abelian category,\allowbreak{} with a finite filtration on every term. For each m fix integers a\_\allowbreak{}m<\allowbreak{}=b\_\allowbreak{}m satisfying F\textasciicircum{}\{a\_\allowbreak{}m\}C\textasciicircum{}m\allowbreak{}=C\textasciicircum{}m and F\textasciicircum{}\{b\_\allowbreak{}m\allowbreak{}+1\}C\textasciicircum{}m\allowbreak{}=0. Zero terms may be given any such bounds. No lower or upper bound on the complex itself is required for this calculation.

For r>\allowbreak{}=1 use the representatives \[
 \mathcal Z_r^{p,m}
     =F^pC^m\cap(d^m)^{-1}(F^{p+r}C^{m+1}).
                                                               \tag{5.1}
\] The source\textquotesingle s quotient formula is equivalently \[
 E_r^{p,m-p}=
 \frac{\mathcal Z_r^{p,m}}
 {\mathcal Z_{r-1}^{p+1,m}
       +d^{m-1}\mathcal Z_{r-1}^{p-r+1,m-1}}.                 \tag{5.2}
\] Both terms in its denominator lie in the numerator:\allowbreak{} the first by its differential condition,\allowbreak{} and the second because it lies in F\^{}p and is a boundary. This formula is obtained from the source\textquotesingle s Z\_\allowbreak{}r and B\_\allowbreak{}r inside F\textasciicircum{}p/\allowbreak{}F\textasciicircum{}\{p\allowbreak{}+1\} by taking representatives and then intersecting the added F\textasciicircum{}\{p\allowbreak{}+1\} with the numerator. No quotient component is discarded.

If \[
 r\ge R(p,m):=
 \max\{1,\ b_{m+1}-p+1,\ p-a_{m-1}+1\},                    \tag{5.3}
\] then F\textasciicircum{}\{p\allowbreak{}+r\}C\textasciicircum{}\{m\allowbreak{}+1\}\allowbreak{}=0 and F\textasciicircum{}\{p-r\allowbreak{}+1\}C\textasciicircum{}\{m-1\}\allowbreak{}=C\textasciicircum{}\{m-1\}. Write Z\textasciicircum{}m\allowbreak{}=ker d\^{}m and B\textasciicircum{}m\allowbreak{}=im d\^{}\{m-1\}. The numerator and denominator in (5.2)\allowbreak{} become exactly \[
 E_r^{p,m-p}=
 \frac{F^pC^m\cap Z^m}
 {(F^{p+1}C^m\cap Z^m)+(F^pC^m\cap B^m)}.                  \tag{5.4}
\] For the second denominator term,\allowbreak{} the image of the inverse image of F\^{}p under d\^{}\{m-1\} is F\^{}p intersect B\^{}m. This identity follows by pulling back the epimorphism C\textasciicircum{}\{m-1\}\allowbreak{}->B\textasciicircum{}m along that intersection,\allowbreak{} so is valid in any abelian category.

The induced filtration on H\textasciicircum{}m(C)\allowbreak{}\allowbreak{}=Z\textasciicircum{}m/\allowbreak{}B\textasciicircum{}m is the image of F\^{}pC\^{}m intersect Z\^{}m. Its associated graded object is precisely (5.4)\allowbreak{},\allowbreak{} by the subobject-quotient isomorphisms. The numerator and both denominator terms remain identical for every larger r. Thus (5.3)\allowbreak{} is an explicit stabilization page,\allowbreak{} with the canonical identification E\_\allowbreak{}r\textasciicircum{}\{p,\allowbreak{}m-p\}\allowbreak{}=gr\textasciicircum{}p H\textasciicircum{}m(C)\allowbreak{}. It proves the convergence comparison itself,\allowbreak{} not only equality of dimensions.

Only a\_\allowbreak{}m<\allowbreak{}=p<\allowbreak{}=b\_\allowbreak{}m can contribute in degree m. Taking a maximum of (5.3)\allowbreak{} over that finite interval gives the single sufficient page \[
 R_m=\max\{1,\ b_{m+1}-a_m+1,\
                   b_m-a_{m-1}+1\}.                         \tag{5.5}
\] The filtration on H\textasciicircum{}m(C)\allowbreak{} has at most b\_\allowbreak{}m-a\_\allowbreak{}m\allowbreak{}+1 potentially nonzero graded pieces:\allowbreak{} it is all of H\^{}m at a\_\allowbreak{}m and zero at b\_\allowbreak{}m\allowbreak{}+1,\allowbreak{} as follows directly from its cycle representatives. It is consequently exhaustive,\allowbreak{} separated and complete; the inverse system of its quotients is eventually H\^{}m itself. Together with (5.3)\allowbreak{} this proves convergence in the source\textquotesingle s exact sense at Homology 6406–\allowbreak{}6421. Boundedness there means finitely many E\_\allowbreak{}0 terms on each total-degree diagonal,\allowbreak{} which is also given by [a\_\allowbreak{}m,\allowbreak{}b\_\allowbreak{}m]\allowbreak{}.

If every term has a common interval [a,\allowbreak{}b]\allowbreak{},\allowbreak{} all degrees stabilize by page b-a\allowbreak{}+1 and each cohomology filtration has at most b-a\allowbreak{}+1 graded pieces. No complex-degree bound is needed for this uniform-window assertion. For varying intervals,\allowbreak{} the adjacent degrees m-1 and m\allowbreak{}+1 in (5.3)\allowbreak{} and (5.5)\allowbreak{} are essential to the calculation and have been retained.

For the filtered derived sequence in section 4,\allowbreak{} PART11 supplies a resolution zero below N whose term in degree m>\allowbreak{}=N has interval \[
 A_m=\min_{N\le j\le m}a_j,\qquad
 B_m=\max_{N\le j\le m}b_j
                                                               \tag{5.6}
\] when [a\_\allowbreak{}j,\allowbreak{}b\_\allowbreak{}j]\allowbreak{} are the input term intervals. T\_\allowbreak{}ext preserves those intervals. Apply (5.3)\allowbreak{} and (5.5)\allowbreak{} with A\_\allowbreak{}m,\allowbreak{}B\_\allowbreak{}m. For the zero terms below N choose A\_\allowbreak{}m\allowbreak{}=A\_\allowbreak{}N and B\_\allowbreak{}m\allowbreak{}=B\_\allowbreak{}N,\allowbreak{} which are valid bounds on zero. This gives an explicit page and filtration length from the input data,\allowbreak{} with their actual integer indices intact.

\subsection{\texorpdfstring{6. The two Cartan–\allowbreak{}Eilenberg sequences and their page shift}{6. The two  sequences and their page shift}}\label{proofs-8352-8772-the-two-cartaneilenberg-sequences-and-their-page-shift}

Original 8558–\allowbreak{}8617 uses the stupid filtration S\^{}pK \allowbreak{}=sigma\_\allowbreak{}\{>\allowbreak{}=p\}K and the canonical filtration G\textasciicircum{}sK\allowbreak{}=tau\_\allowbreak{}\{<\allowbreak{}=-s\}K. The complete comparison was proved in PART10 section 3,\allowbreak{} including the injective filtrations on the actual total complex and the formula \[
 E_r^{s,t}(G)=
 {}''E_{r+1}^{2s+t,-s}\quad(r\ge1).                          \tag{6.1}
\] The new review in sections 2–\allowbreak{}4 confirms that the filtered derived functor used there has the precise associated-graded and forgetful comparisons (4.1)\allowbreak{}. For S,\allowbreak{} the graded piece is K\textasciicircum{}p[-p]\allowbreak{},\allowbreak{} so its derived E\_\allowbreak{}1 is R\textasciicircum{}qT(K\textasciicircum{}p)\allowbreak{}. For G,\allowbreak{} the graded piece is quasi-isomorphic to H\textasciicircum{}\{-s\}(K)\allowbreak{}[s]\allowbreak{},\allowbreak{} so its derived E\_\allowbreak{}1 is R\textasciicircum{}\{2s\allowbreak{}+t\}T(H\textasciicircum{}\{-s\}(K)\allowbreak{})\allowbreak{}. The exact cycle and boundary quotient calculation establishing (6.1)\allowbreak{} retains the original total differential d\_\allowbreak{}1\allowbreak{}+(-1)\allowbreak{}\textasciicircum{}p d\_\allowbreak{}2. It identifies the differentials as well:\allowbreak{} bidegree (r,\allowbreak{}1-r)\allowbreak{} becomes (r\allowbreak{}+1,\allowbreak{}-r)\allowbreak{}. Thus the source\textquotesingle s two promised sequences agree from their canonical E\_\allowbreak{}1 and E\_\allowbreak{}2 pages,\allowbreak{} respectively.

The split horizontal kernel sequences in the Cartan–\allowbreak{}Eilenberg resolution make this argument valid for additive T as proved in PART10. We do not assume that arbitrary underived kernels commute with T. RECON-066 /\allowbreak{} P02-E0065 /\allowbreak{} DERIVED-059 is the existing MC-STK-ERR-0870 insertion of ``is''. RECON-067 /\allowbreak{} DERIVED-060 is MC-STK-ERR-0871,\allowbreak{} the existing idiom correction at 8603.

\subsection{7. Ext groups and their exact boundary maps}\label{proofs-8352-8772-ext-groups-and-their-exact-boundary-maps}

Keep the source\textquotesingle s definition Ext\textasciicircum{}i\_\allowbreak{}A(X,\allowbreak{}Y)\allowbreak{}\allowbreak{}=Hom\_\allowbreak{}D(A)\allowbreak{}(X,\allowbreak{}Y[i]\allowbreak{})\allowbreak{}. The equality with Hom\_\allowbreak{}D(A)\allowbreak{}(X[-i]\allowbreak{},\allowbreak{}Y)\allowbreak{} is the natural bijection given by the inverse shift autoequivalence. For a triangle Y\allowbreak{}->Y'\allowbreak{}->Y'' --h-\allowbreak{}->Y[1]\allowbreak{},\allowbreak{} the connecting map on Ext sends f:\allowbreak{}X\allowbreak{}->Y''[i]\allowbreak{} to h[i]\allowbreak{}f:\allowbreak{}X\allowbreak{}->Y[i\allowbreak{}+1]\allowbreak{}. The long exact Hom sequence for that triangle proves exactness with the displayed objects. For a triangle X\allowbreak{}->X'\allowbreak{}->X'' --k-\allowbreak{}->X[1]\allowbreak{},\allowbreak{} the connecting map sends f:\allowbreak{}X\allowbreak{}->Y[i]\allowbreak{} to f[1]\allowbreak{}k:\allowbreak{}X''\allowbreak{}->Y[i\allowbreak{}+1]\allowbreak{}. This is the shifted contravariant Hom boundary. These are the actual maps in 8654–\allowbreak{}8666. Fullness of the bounded derived subcategories allows the same computation when both inputs belong to one of them,\allowbreak{} since their shifts stay there.

If Y\allowbreak{}->I is a bounded-below injective resolution,\allowbreak{} the target I[i]\allowbreak{} remains such an injective complex,\allowbreak{} with lower bound shifted by -i and differential (-1)\allowbreak{}\textasciicircum{}i d\_\allowbreak{}I. The injective Hom comparison identifies Ext\textasciicircum{}i(X,\allowbreak{}Y)\allowbreak{} with Hom\_\allowbreak{}K(A)\allowbreak{}(X,\allowbreak{}I[i]\allowbreak{})\allowbreak{}. The source X need not have any term bound. Dually,\allowbreak{} for a bounded-above projective resolution P\allowbreak{}->X,\allowbreak{} P[-i]\allowbreak{} has upper bound shifted by i and its full shift differential. The projective Hom comparison gives Ext\textasciicircum{}i(X,\allowbreak{}Y)\allowbreak{}\allowbreak{}= Hom\_\allowbreak{}K(A)\allowbreak{}(P[-i]\allowbreak{},\allowbreak{}Y)\allowbreak{} without a bound on Y. This proves 8681–\allowbreak{}8706 with the source\textquotesingle s opposite bounds.

\subsection{8. The exact lowest Ext degree}\label{proofs-8352-8772-the-exact-lowest-ext-degree}

Assume H\textasciicircum{}j(X)\allowbreak{}\allowbreak{}=0 for j\textgreater a and H\textasciicircum{}j(Y)\allowbreak{}\allowbreak{}=0 for j\textless b. Use the original canonical truncations:\allowbreak{} Y\allowbreak{}->Y\_\allowbreak{}b\allowbreak{}=tau\_\allowbreak{}\{>\allowbreak{}=b\}Y is a quasi-isomorphism; a denominator L\allowbreak{}->X may be replaced by L\_\allowbreak{}a\allowbreak{}=tau\_\allowbreak{}\{<\allowbreak{}=a\}L\allowbreak{}->X. Thus every derived map X\allowbreak{}->Y[n]\allowbreak{} has a roof whose numerator is a chain map L\_\allowbreak{}a\allowbreak{}->Y\_\allowbreak{}b[n]\allowbreak{}. Here L\_\allowbreak{}a\textasciicircum{}j\allowbreak{}=0 for j>a,\allowbreak{} and Y\_\allowbreak{}b\textasciicircum{}j\allowbreak{}=0 for j\textless b. For n\textless b-a and j<\allowbreak{}=a,\allowbreak{} j\allowbreak{}+n<b,\allowbreak{} so the target component is zero. For j\textgreater a the source component is zero. All components of the numerator vanish,\allowbreak{} proving Ext\textasciicircum{}n(X,\allowbreak{}Y)\allowbreak{}\allowbreak{}=0.

At n\allowbreak{}=b-a only the component in degree a can be nonzero. The chain equations make it kill im d\_\allowbreak{}\{L\_\allowbreak{}a\}\textasciicircum{}\{a-1\} and land in ker d\_\allowbreak{}\{Y\_\allowbreak{}b\}\textasciicircum{}b. The target shift multiplies that last differential by (-1)\allowbreak{}\textasciicircum{}\{b-a\},\allowbreak{} which has the same kernel; the sign has not been removed from the shifted complex. Therefore the map factors uniquely as \[
 L_a^a\longrightarrow H^a(L_a)
 \xrightarrow{\bar f} H^b(Y_b)
 \longrightarrow Y_b^b.                                    \tag{8.1}
\] There are no nonzero homotopies between such maps:\allowbreak{} a homotopy component in degree j has source L\_\allowbreak{}a\textasciicircum{}j and target Y\_\allowbreak{}b\textasciicircum{}\{j-1\allowbreak{}+b-a\}; for j<\allowbreak{}=a the target is below b,\allowbreak{} and for j\textgreater a the source is zero.

The denominator induces H\textasciicircum{}a(L\_\allowbreak{}a)\allowbreak{}\textasciitilde{}\allowbreak{}=H\textasciicircum{}a(X)\allowbreak{}. Consequently a roof gives the map bar f H\textasciicircum{}a(s)\allowbreak{}\textasciicircum{}\{-1\} from H\textasciicircum{}a(X)\allowbreak{} to H\textasciicircum{}b(Y)\allowbreak{}. Refining the roof preserves that map,\allowbreak{} by the cohomology equation of the refining denominator. Conversely,\allowbreak{} given H\textasciicircum{}a(X)\allowbreak{}\allowbreak{}->H\textasciicircum{}b(Y)\allowbreak{},\allowbreak{} take tau\_\allowbreak{}\{<\allowbreak{}=a\}X,\allowbreak{} map its top term to H\textasciicircum{}a(X)\allowbreak{},\allowbreak{} follow the given map and the cycle inclusion into tau\_\allowbreak{}\{>\allowbreak{}=b\}Y[b-a]\allowbreak{},\allowbreak{} and use zero maps in every other degree. It is a chain map by (8.1)\allowbreak{} and defines a roof with precisely the given cohomology map. A roof and this constructed one agree after a common denominator,\allowbreak{} because their only component is the unique factorization (8.1)\allowbreak{}. This proves the natural isomorphism \[
 \operatorname{Ext}^{b-a}_A(X,Y)
   \xrightarrow{\sim}
 \operatorname{Hom}_A(H^a(X),H^b(Y)).                         \tag{8.2}
\] No injectives,\allowbreak{} projectives,\allowbreak{} or replacement objects with altered support are assumed.

RECON-068 /\allowbreak{} P02-E0067 /\allowbreak{} DERIVED-064 retains MC-STK-ERR-0875:\allowbreak{} after tau\_\allowbreak{}\{<\allowbreak{}=a\},\allowbreak{} the vanishing is L\textasciicircum{}j\allowbreak{}=0 for j>a,\allowbreak{} not j\textless a. The report\textquotesingle s mention of cohomology is supporting intuition; the corrected printed statement is about the terms,\allowbreak{} exactly as needed in this proof. The case a\allowbreak{}=b\allowbreak{}=0 gives Ext\textasciicircum{}n(B,\allowbreak{}A)\allowbreak{}\allowbreak{}=0 for n\textless0 and Ext\textasciicircum{}0(B,\allowbreak{}A)\allowbreak{}\allowbreak{}=Hom(B,\allowbreak{}A)\allowbreak{}.

Finally a short exact sequence of objects gives its canonical triangle in D(A)\allowbreak{}. Applying the two Hom sequences of section 7 and this negative-degree vanishing produces the two displayed long exact sequences at 8753–\allowbreak{}8765,\allowbreak{} including their initial zero terms. They therefore exist without enough injectives or enough projectives. The Yoneda construction following 8772 is the next substantive review,\allowbreak{} not assumed as proof of the assertions above.

\subsection{9. Receiving use and the next review}\label{proofs-8352-8772-receiving-use-and-the-next-review}

For the bounded-below ringed-space construction in current Cohomology 7307–\allowbreak{}7339,\allowbreak{} take its actual injective resolution I and T\allowbreak{}=Gamma(X,\allowbreak{}-)\allowbreak{}. The split filtration on each I\^{}m makes Gamma(X,\allowbreak{}F\textasciicircum{}p I\textasciicircum{}m)\allowbreak{} the subobject specified in (3.1)\allowbreak{}. If the input intervals are [a\_\allowbreak{}j,\allowbreak{}b\_\allowbreak{}j]\allowbreak{},\allowbreak{} PART11\textquotesingle s construction and (5.6)\allowbreak{} give explicit intervals [A\_\allowbreak{}m,\allowbreak{}B\_\allowbreak{}m]\allowbreak{} on Gamma(X,\allowbreak{}I\textasciicircum{}m)\allowbreak{}. Equation (5.5)\allowbreak{} gives its stabilization page,\allowbreak{} and (5.4)\allowbreak{} identifies its graded abutment. For a common input interval [a,\allowbreak{}b]\allowbreak{} the page is b-a\allowbreak{}+1. This strengthens the receiver\textquotesingle s qualitative finite- filtration conclusion. Its unbounded K-injective variant at 7286–\allowbreak{}7305 has separate hypotheses; this review does not assert that its filtration intervals satisfy the bounded-below construction (5.6)\allowbreak{}.

The original continuation through 8930 was read but its Yoneda proofs are not yet adjudicated. In particular,\allowbreak{} the extension-composition display at 8913–\allowbreak{}8929 appears to reverse the endpoint variance relative to the following extensions. The next review must check the whole composition passage and its sign and roof conventions before recording an operation. This is a routed observation,\allowbreak{} not a counted correction or an accepted Yoneda comparison.
